% GENERATED FILE. Edit canonical lessons, then run npm run generate:books.
% canonical-source-hash: fnv1a64:d3a98f20324e70c9
% canonical-lessons: GU-C268-ladvu-fight, GU-C268-vadhvu-scold, GU-C268-chetarvu, GU-C268-manavvu, GU-C268-salah

\chapter{Quarrels and advice}
\label{ch:gu-quarrels-and-advice}
\hlchaptermodality{268}

\begin{chapteropening}
\textbf{By the end of this chapter:} \emph{I can say that people do not quarrel and ask who cheated or persuaded whom.}

Complete the last lesson of chapter 268: I can say that people do not quarrel and ask who cheated or persuaded whom.
\end{chapteropening}

\section[\texorpdfstring{\gu{લડવું} (laḍvũ)}{લડવું (laḍvũ)}]{\gu{લડવું} (laḍvũ) — to quarrel, to fight}

\label{lesson:GU-C268-ladvu-fight}



\begin{quote}
Before the new one: say the Gujarati for to get irritated, then the Gujarati for to sulk.
\end{quote}

\subsection*{You'll want to know: \gu{લડવું}}
\textbf{\gu{લડવું}} — \emph{laḍvũ} — \textquotedblleft{}to quarrel, to fight\textquotedblright{}.

\begin{quote}
\textbf{\gu{અમે} \gu{ક્યારેય} \gu{નથી} \gu{લડતા}.} — \emph{ame kyārey nathī laḍtā} — \textquotedblleft{}We never quarrel.\textquotedblright{}
\end{quote}

\begin{cousinweb}
\textbf{\gu{લડવું}} (\emph{laḍvũ}) is a tadbhava. It comes from \emph{*laḍ}, \textquotedblleft{}fight\textquotedblright{}, a reconstructed form; CDIAL glosses the Gujarati verb \textquotedblleft{}to fight\textquotedblright{} (CDIAL 10920).
\end{cousinweb}

\begin{grammarlens}[title={What you've built}]
One more word for everyday life at A2.
\end{grammarlens}

\subsection*{Your turn}
\begin{itemize}
\raggedright
  \item \emph{Say it:} \emph{laḍvũ}
  \item \emph{Say it:} \emph{laḍvũ}, once more
  \item \emph{Say it:} \emph{ame kyārey nathī laḍtā} — We never quarrel.
  \item \emph{Recall:} say the Gujarati for to get irritated, then the Gujarati for to sulk, then say \emph{laḍvũ} again
\end{itemize}

\begin{tcolorbox}[breakable,colback=teal!4,colframe=teal!35!black,title={Before you move on}]
What does \gu{લડવું} mean? (\textquotedblleft{}To quarrel, to fight\textquotedblright{}.) Say it once more.
\end{tcolorbox}
\section[\texorpdfstring{\gu{વઢવું} (vaḍhvũ)}{વઢવું (vaḍhvũ)}]{\gu{વઢવું} (vaḍhvũ) — to scold}

\label{lesson:GU-C268-vadhvu-scold}



\begin{quote}
Before the new one: say the Gujarati for to sulk, then the Gujarati for to quarrel.
\end{quote}

\subsection*{You'll want to know: \gu{વઢવું}}
\textbf{\gu{વઢવું}} — \emph{vaḍhvũ} — \textquotedblleft{}to scold\textquotedblright{}.

\begin{quote}
\textbf{\gu{મારી} \gu{માતા} \gu{મને} \gu{નથી} \gu{વઢતી}.} — \emph{mārī mātā mane nathī vaḍhtī} — \textquotedblleft{}My mother does not scold me.\textquotedblright{}
\end{quote}

\begin{cousinweb}
The origin of \textbf{\gu{વઢવું}} (\emph{vaḍhvũ}) is not settled. CDIAL lists no Gujarati verb of this form with the meaning \textquotedblleft{}to scold\textquotedblright{}, and our other sources do not trace it.
\end{cousinweb}

\begin{grammarlens}[title={What you've built}]
One more word for everyday life at A2.
\end{grammarlens}

\subsection*{Your turn}
\begin{itemize}
\raggedright
  \item \emph{Say it:} \emph{vaḍhvũ}
  \item \emph{Say it:} \emph{vaḍhvũ}, once more
  \item \emph{Say it:} \emph{mārī mātā mane nathī vaḍhtī} — My mother does not scold me.
  \item \emph{Recall:} say the Gujarati for to sulk, then the Gujarati for to quarrel, then say \emph{vaḍhvũ} again
\end{itemize}

\begin{tcolorbox}[breakable,colback=teal!4,colframe=teal!35!black,title={Before you move on}]
What does \gu{વઢવું} mean? (\textquotedblleft{}To scold\textquotedblright{}.) Say it once more.
\end{tcolorbox}
\section[\texorpdfstring{\gu{છેતરવું} (chetarvũ)}{છેતરવું (chetarvũ)}]{\gu{છેતરવું} (chetarvũ) — to deceive, to cheat}

\label{lesson:GU-C268-chetarvu}



\begin{quote}
Before the new one: say the Gujarati for to quarrel, then the Gujarati for to scold.
\end{quote}

\subsection*{You'll want to know: \gu{છેતરવું}}
\textbf{\gu{છેતરવું}} — \emph{chetarvũ} — \textquotedblleft{}to deceive, to cheat\textquotedblright{}.

\begin{quote}
\textbf{\gu{શું} \gu{તેણે} \gu{તમને} \gu{છેતર્યા}?} — \emph{śũ teṇe tamne chetaryā?} — \textquotedblleft{}Did he cheat you?\textquotedblright{}
\end{quote}

\begin{cousinweb}
The origin of \textbf{\gu{છેતરવું}} (\emph{chetarvũ}) is not settled. CDIAL lists no Gujarati verb of this shape; its Gujarati \emph{chaḷvũ}, \textquotedblleft{}to deceive\textquotedblright{} (CDIAL 5003), is a different word.
\end{cousinweb}

\begin{grammarlens}[title={What you've built}]
One more word for everyday life at A2.
\end{grammarlens}

\subsection*{Your turn}
\begin{itemize}
\raggedright
  \item \emph{Say it:} \emph{chetarvũ}
  \item \emph{Say it:} \emph{chetarvũ}, once more
  \item \emph{Say it:} \emph{śũ teṇe tamne chetaryā?} — Did he cheat you?
  \item \emph{Recall:} say the Gujarati for to quarrel, then the Gujarati for to scold, then say \emph{chetarvũ} again
\end{itemize}

\begin{tcolorbox}[breakable,colback=teal!4,colframe=teal!35!black,title={Before you move on}]
What does \gu{છેતરવું} mean? (\textquotedblleft{}To deceive, to cheat\textquotedblright{}.) Say it once more.
\end{tcolorbox}
\section[\texorpdfstring{\gu{મનાવવું} (manāvvũ)}{મનાવવું (manāvvũ)}]{\gu{મનાવવું} (manāvvũ) — to persuade, to coax}

\label{lesson:GU-C268-manavvu}



\begin{quote}
Before the new one: say the Gujarati for to scold, then the Gujarati for to deceive.
\end{quote}

\subsection*{You'll want to know: \gu{મનાવવું}}
\textbf{\gu{મનાવવું}} — \emph{manāvvũ} — \textquotedblleft{}to persuade, to coax\textquotedblright{}.

\begin{quote}
\textbf{\gu{શું} \gu{તમે} \gu{તેને} \gu{મનાવશો}?} — \emph{śũ tame tene manāvśo?} — \textquotedblleft{}Will you persuade him?\textquotedblright{}
\end{quote}

The causative of \textbf{\gu{માનવું}}, \emph{mānvũ}, \textquotedblleft{}to believe\textquotedblright{}: to bring someone round.

\begin{cousinweb}
\textbf{\gu{મનાવવું}} (\emph{manāvvũ}) is a tadbhava. CDIAL lists it, glossed \textquotedblleft{}to persuade\textquotedblright{}, beside \emph{mānvũ} under Sanskrit \emph{mányatē}, \textquotedblleft{}thinks, honours; agrees\textquotedblright{} (CDIAL 9857).
\end{cousinweb}

\begin{grammarlens}[title={What you've built}]
One more word for everyday life at A2.
\end{grammarlens}

\subsection*{Your turn}
\begin{itemize}
\raggedright
  \item \emph{Say it:} \emph{manāvvũ}
  \item \emph{Say it:} \emph{manāvvũ}, once more
  \item \emph{Say it:} \emph{śũ tame tene manāvśo?} — Will you persuade him?
  \item \emph{Recall:} say the Gujarati for to scold, then the Gujarati for to deceive, then say \emph{manāvvũ} again
\end{itemize}

\begin{tcolorbox}[breakable,colback=teal!4,colframe=teal!35!black,title={Before you move on}]
What does \gu{મનાવવું} mean? (\textquotedblleft{}To persuade, to coax\textquotedblright{}.) Say it once more.
\end{tcolorbox}
\section[\texorpdfstring{\gu{સલાહ} (salāh)}{સલાહ (salāh)}]{\gu{સલાહ} (salāh) — advice}

\label{lesson:GU-C268-salah}



\begin{quote}
Before the new one: say the Gujarati for to deceive, then the Gujarati for to persuade.
\end{quote}

\subsection*{You'll want to know: \gu{સલાહ}}
\textbf{\gu{સલાહ}} — \emph{salāh} — \textquotedblleft{}advice\textquotedblright{}.

\begin{quote}
\textbf{\gu{તે} \gu{કોઈની} \gu{સલાહ} \gu{નથી} \gu{માનતો}.} — \emph{te koīnī salāh nathī mānto} — \textquotedblleft{}He does not take anyone's advice.\textquotedblright{}
\end{quote}

A feminine noun.

\begin{cousinweb}
\textbf{\gu{સલાહ}} (\emph{salāh}) is a loan usually traced to Arabic \emph{ṣalāḥ}, \textquotedblleft{}soundness, goodness\textquotedblright{}, which reached Gujarati through Persian. Its Arabic root is \emph{ṣ-l-ḥ}, \textquotedblleft{}to be good, to be put right\textquotedblright{} (Lane).
\end{cousinweb}

\begin{grammarlens}[title={What you've built}]
One more word for everyday life at A2.
\end{grammarlens}

\subsection*{Your turn}
\begin{itemize}
\raggedright
  \item \emph{Say it:} \emph{salāh}
  \item \emph{Say it:} \emph{salāh}, once more
  \item \emph{Say it:} \emph{te koīnī salāh nathī mānto} — He does not take anyone's advice.
  \item \emph{Recall:} say the Gujarati for to deceive, then the Gujarati for to persuade, then say \emph{salāh} again
\end{itemize}

\begin{tcolorbox}[breakable,colback=teal!4,colframe=teal!35!black,title={Before you move on}]
What does \gu{સલાહ} mean? (\textquotedblleft{}Advice\textquotedblright{}.) Say it once more.
\end{tcolorbox}
